\documentclass[10pt,a4paper]{article}

\usepackage[T1]{fontenc}
\usepackage[utf8]{inputenc}
\usepackage[ngerman]{babel}
\usepackage{lmodern}
\usepackage{geometry}
\geometry{left=1.75cm,right=1.75cm,top=1.55cm,bottom=1.55cm}
\usepackage{setspace}
\setstretch{0.98}
\usepackage{enumitem}
\usepackage{hyperref}
\hypersetup{colorlinks=true,linkcolor=blue,urlcolor=blue}
\setlist[itemize]{leftmargin=1.1em,itemsep=0.04em,topsep=0.04em,parsep=0pt,partopsep=0pt}
\pagestyle{empty}

\begin{document}

\begin{center}
{\Large \textbf{Founder-One-Pager: Aktueller Produktstand}}\\[0.15em]
{\large Triagent: Evidence-first Phishing-Triage für regulierte Organisationen}\\[0.35em]
\textbf{Alexander Hülsmann} -- Masterstudent Informatik, Universität Münster
\end{center}
\small

\vspace{0.25em}
\noindent\textbf{Kurzfassung}\\
Triagent ist inzwischen kein reines Validierungskonzept mehr, sondern ein \textbf{lauffähiger On-Prem-Phishing-Triage-Demo-Stack}. Die Kernidee bleibt: Analysten sollen Phishing-Meldungen \textbf{nicht über verstreute Einzelschritte und Einzeltools} prüfen, sondern in einer Oberfläche technische Evidenz, Entscheidung und Dokumentation zusammenführen. Das Produkt ist heute \textbf{evidence-first, analyst-in-the-loop und kampagnenfähig}; es ist \textbf{noch nicht enterprise-fertig}, aber die Kern-Workflows sind als funktionierender Prototyp vorhanden.

\vspace{0.3em}
\noindent\textbf{Was heute konkret gebaut ist}
\begin{itemize}
  \item \textbf{Analyst Workspace:} Web-App mit Uploads, In-Tray, Dashboard, Admin-Bereich und vollständiger Report-Ansicht.
  \item \textbf{E-Mail-Intake:} Upload und Parsing von \texttt{.eml} und \texttt{.msg}, plus Outlook-Add-in für Meldungen aus dem Client.
  \item \textbf{Technische Analyse:} Header-/Auth-Auswertung (SPF, DKIM, DMARC), URL-Extraktion mit Redirect-Auflösung, Attachment-Metadaten, Transmission-Hop-Ansicht, Raw Source/X-Headers.
  \item \textbf{Analystische Hilfen:} Same-org Domain-Lookalike-Detection, ATT\&CK-Mapping, IOC-Export, Risiko-/Verdachtskontext und Flagging einzelner Artefakte.
  \item \textbf{Evidenz \& Nachvollziehbarkeit:} Export als JSON-Bundle, Markdown und PDF; zusätzlich IOC-Export als JSON/CSV; Original-Mail und Anhänge bleiben als Artefakte erhalten und downloadbar.
  \item \textbf{Betrieb \& Governance:} RBAC, API Keys, Audit Log, Docker-Compose-Deployment sowie synthetischer Testkorpus für reproduzierbare Demos und Regressionstests.
\end{itemize}

\noindent\textbf{Für wen}
\begin{itemize}
  \item \textbf{Nutzer:} SOC-Analysten / Phishing-Triage-Teams
  \item \textbf{Champion:} Analyst oder Teamlead mit hohem Pain bei Mail-Wellen und Audit-/Doku-Aufwand
  \item \textbf{Käufer:} SOC-Leitung / CISO in Umfeldern mit Datenhoheit, Compliance- und On-Prem-Anforderungen
\end{itemize}

\noindent\textbf{Welches Problem Triagent adressiert}
\begin{itemize}
  \item Triage ist oft \textbf{tool-fragmentiert}: Header/Auth prüfen, URLs öffnen, Anhänge sichern, Doku schreiben, IOCs weiterreichen.
  \item In regulierten Umfeldern sind \textbf{externe Uploads oder SaaS-Workflows} oft politisch oder technisch schwierig.
  \item Viel Zeit geht nicht nur in die Entscheidung, sondern in \textbf{Begründung, Export und Auditierbarkeit}.
  \item Bei Phishing-Wellen entsteht zusätzlich \textbf{redundante Mehrfachprüfung ähnlicher Meldungen}.
\end{itemize}

\noindent\textbf{Warum ich glaube, dass der Ansatz interessant ist}
\begin{itemize}
  \item \textbf{Kein generischer ``AI SOC Copilot'':} enger, klarer Workflow mit überprüfbarer Evidenz.
  \item \textbf{On-Prem + nachvollziehbar:} Daten bleiben lokal, Original-Mail bleibt erhalten, Entscheidungen sind exportierbar.
  \item \textbf{Nicht nur Verdict, sondern Investigation Output:} Analysten bekommen Artefakte, IOC-Export, ATT\&CK-Kontext und Bericht in einem Lauf.
  \item \textbf{Campaign Layer ist angelegt:} Kampagnen-Clustering, Zuordnung und Maintenance existieren im Backend/Data Model bereits; die eigentliche Produktvalidierung dafür ist der nächste große Schritt.
\end{itemize}

\noindent\textbf{Was noch offen ist}
\begin{itemize}
  \item Heute ist es ein \textbf{starker Demo-/Pilot-Prototyp}, aber noch kein produktionsreifes Enterprise-Produkt.
  \item Es fehlen noch Dinge wie \textbf{SSO/IdP-Integration, tiefe Ticketing-/SIEM-Integrationen, härtere Packaging-/Ops-Story und echte Pilot-Feedbackschleifen}.
  \item Die größte offene Produktfrage ist: \textbf{Ist die erste Keilfunktion das einzelne Evidence-first-Report-Handling oder das kampagnenweise Arbeiten?}
\end{itemize}

\noindent\textbf{Wobei ein Co-Founder besonders helfen könnte}
\begin{itemize}
  \item Produkt- und GTM-Schärfung: \textbf{welcher Entry Point zuerst wirklich zieht}
  \item Design-Partner- und Pilotzugang in Security-/SOC-Umfeldern
  \item Ergänzende Stärke in \textbf{Security-Produktisierung, Enterprise-Engineering oder kommerziellem Aufbau}
\end{itemize}

\end{document}
